\section{Problem 2: Disk Pendulum}
\begin{figure} [h]
\centering
\includegraphics[width=0.5\textwidth]{figures/Problem_2.png}
\caption{Disk pendulum system.}
\end{figure}
We use the same generalized coordinates and reference frames from HW01:
\[
q_1 = \phi(t),
\qquad
q_2 = \theta(t)
\]
where:
\begin{itemize}
    \item $\phi$ is the rotation angle of the vertical shaft,
    \item $\theta$ is the pendulum swing angle measured from the downward vertical.
\end{itemize}

From HW01, the velocity of the point mass is:
\[
\mathbf{v}_{m_b}
=
\left(R_d\dot\phi + L\dot\phi\cos\theta\right)\mathbf{b}_2
-
L\dot\theta\mathbf{b}_3
\]

Hence:
\[
|\mathbf{v}_{m_b}|^2
=
\left(R_d\dot\phi + L\dot\phi\cos\theta\right)^2
+
L^2\dot\theta^2
\]

\subsection*{(a) Total Kinetic Energy}

The total kinetic energy consists of:
\begin{itemize}
    \item rotational kinetic energy of the disk,
    \item translational kinetic energy of the point mass.
\end{itemize}

The disk kinetic energy is:
\[
T_d = \frac12 I_d \dot\phi^2
\]

The point-mass kinetic energy is:
\begin{align}
T_b
&= \frac12 m_b |\mathbf{v}_{m_b}|^2
\\
&=
\frac12 m_b
\left[
\left(R_d\dot\phi + L\dot\phi\cos\theta\right)^2
+
L^2\dot\theta^2
\right]
\end{align}

Therefore the total kinetic energy is:
\begin{align}
T
&= T_d + T_b
\\
&=
\frac12 I_d\dot\phi^2
+
\frac12 m_b
\left(R_d + L\cos\theta\right)^2 \dot\phi^2
+
\frac12 m_bL^2\dot\theta^2
\end{align}

\subsection*{(b) Potential Energy}

Choose the zero potential energy datum at the pivot point $P$.

The vertical position of the point mass is:
\[
z = -L\cos\theta
\]

Hence the potential energy is:
\[
V = m_bgL\sin\theta
\]

\subsection*{(c) Lagrange Equations}

The Lagrangian is:
\[
\mathcal{L} = T - V
\]

Substituting the expressions for $T$ and $V$:
\begin{align}
\mathcal{L}
&=
\frac12 I_d\dot\phi^2
+
\frac12 m_b
\left(R_d + L\cos\theta\right)^2\dot\phi^2
\\
&\quad
+
\frac12 m_bL^2\dot\theta^2
-
m_bgL\sin\theta
\end{align}

\subsubsection*{Equation for $\phi$}

Compute:
\[
\frac{\partial \mathcal{L}}{\partial \dot\phi}
=
\left[
I_d + m_b\left(R_d+L\cos\theta\right)^2
\right]\dot\phi
\]

Taking the time derivative:
\begin{align}
\frac{d}{dt}
\left(
\frac{\partial \mathcal{L}}{\partial \dot\phi}
\right)
&=
\left[
I_d + m_b(R_d+L\cos\theta)^2
\right]\ddot\phi
\\
&\quad
-2m_bL(R_d+L\cos\theta)\sin\theta\dot\theta\dot\phi
\end{align}

Since:
\[
\frac{\partial \mathcal{L}}{\partial \phi} = 0
\]

The equation of motion for $\phi$ is:
\begin{align}
&\left[
I_d + m_b(R_d+L\cos\theta)^2
\right]\ddot\phi
\\
&\quad
-2m_bL(R_d+L\cos\theta)\sin\theta\dot\theta\dot\phi
=0
\end{align}

\subsubsection*{Equation for $\theta$}

Compute:
\[
\frac{\partial \mathcal{L}}{\partial \dot\theta}
= m_bL^2\dot\theta
\]

Thus:
\[
\frac{d}{dt}
\left(
\frac{\partial \mathcal{L}}{\partial \dot\theta}
\right)
= m_bL^2\ddot\theta
\]

Now compute:
\begin{align}
\frac{\partial \mathcal{L}}{\partial \theta}
&=
-m_bL(R_d+L\cos\theta)\sin\theta\dot\phi^2
\\
&\quad
-m_bgL\cos\theta
\end{align}

Hence the equation of motion for $\theta$ is:
\[
m_bL^2\ddot\theta
+
m_bL(R_d+L\cos\theta)\sin\theta\dot\phi^2
-
m_bgL\cos\theta
=0
\]

\subsection*{(d) Matrix Form}

Define:
\[
\mathbf{q}
=
\begin{bmatrix}
\phi \\
\theta
\end{bmatrix}
\]

The equations of motion can be written as:
\[
\mathbf{M}(q)\ddot q
+
\mathbf{h}(q,\dot q)
+
\mathbf{g}(q)
=0
\]

The mass matrix is:
\[
\mathbf{M}(q)
=
\begin{bmatrix}
I_d + m_b(R_d+L\cos\theta)^2 & 0 \\
0 & m_bL^2
\end{bmatrix}
\]

The velocity-coupling vector is:
\[
\mathbf{h}(q,\dot q)
=
\begin{bmatrix}
-2m_bL(R_d+L\cos\theta)\sin\theta\dot\theta\dot\phi \\
m_bL(R_d+L\cos\theta)\sin\theta\dot\phi^2
\end{bmatrix}
\]

The gravity vector is:
\[
\mathbf{g}(q)
=
\begin{bmatrix}
0 \\
-m_bgL\cos\theta
\end{bmatrix}
\]

\subsection*{(e) Cyclic Coordinate}

Since the Lagrangian does not explicitly depend on $\phi$:
\[
\frac{\partial \mathcal{L}}{\partial \phi}=0
\]

Therefore $\phi$ is a cyclic (ignorable) coordinate.

The conserved generalized momentum is:
\begin{align}
p_\phi
&=
\frac{\partial \mathcal{L}}{\partial \dot\phi}
\\
&=
\left[
I_d + m_b(R_d+L\cos\theta)^2
\right]\dot\phi
\end{align}

This quantity represents the total angular momentum about the vertical shaft axis.

\subsection*{(f) Virtual Work}

Suppose an external horizontal force $\mathbf{F}$ acts at the point mass $m_b$.

The position vector of the point mass is:
\[
\mathbf{r}_{m_b}
=
R_d\mathbf{a}_1 + L\mathbf{b}_1
\]

The virtual displacement is:
\[
\delta \mathbf{r}_{m_b}
=
\frac{\partial \mathbf{r}_{m_b}}{\partial \phi}\delta\phi
+
\frac{\partial \mathbf{r}_{m_b}}{\partial \theta}\delta\theta
\]

The virtual work is:
\[
\delta W
=
\mathbf{F}\cdot \delta\mathbf{r}_{m_b}
\]

Comparing with:
\[
\delta W = Q_\phi\delta\phi + Q_\theta\delta\theta
\]

The generalized forces are:
\begin{align}
Q_\phi
&=
\mathbf{F}\cdot
\frac{\partial \mathbf{r}_{m_b}}{\partial \phi}
\\
Q_\theta
&=
\mathbf{F}\cdot
\frac{\partial \mathbf{r}_{m_b}}{\partial \theta}
\end{align}

\subsection*{(g) Conservation of Mechanical Energy}

For the unforced and frictionless system:
\[
E = T + V
\]

Substituting the expressions:
\begin{align}
E
&=
\frac12 I_d\dot\phi^2
+
\frac12 m_b(R_d+L\cos\theta)^2\dot\phi^2
\\
&\quad
+
\frac12 m_bL^2\dot\theta^2
-
m_bgL\cos\theta
\end{align}

Since there are no non-conservative forces acting on the system:
\[
\frac{dE}{dt}=0
\]

Therefore the total mechanical energy is conserved.

\subsection*{(h) Bonus: Linearization About Steady Spin}

Consider small oscillations about:
\[
\theta = 0,
\qquad
\dot\phi = \Omega_0
\]

Using the approximations:
\[
\sin\theta \approx \theta,
\qquad
\cos\theta \approx 1
\]

The $\theta$ equation becomes:
\begin{align}
m_bL^2\ddot\theta
&+
m_bL(R_d+L)\Omega_0^2\theta
\\
&+
m_bgL\theta = 0
\end{align}

Dividing by $m_bL^2$:
\[
\ddot\theta
+
\left(
\frac{(R_d+L)\Omega_0^2}{L}
+
\frac{g}{L}
\right)\theta
=0
\]

Hence the linearized small-angle equation is:
\[
\boxed{
\ddot\theta
+
\omega_n^2\theta = 0
}
\]
with:
\[
\omega_n^2
=
\frac{(R_d+L)\Omega_0^2}{L}
+
\frac{g}{L}
\]
